Symmetry is one of the few inductive biases that reliably pays for itself in
machine learning, from group-equivariant convolutions
\citep{cohen2016group} to general constructions for
arbitrary matrix groups \citep{finzi2021practical} and the geometric programme
that organises them \citep{bronstein2021geometric}. All of it presumes that
the symmetry is known. When it is not, it has to be discovered from data, and
the question this paper is about is what a discovery may honestly claim.

A continuous symmetry of a dynamical system $\dot x = F(x)$ is a one-parameter
family of state transformations that leaves the dynamics unchanged. Each such
family is generated by a vector field $\xi$, and $\xi$ generates a symmetry
exactly when the Lie bracket $[\xi,F]$ vanishes identically. Because the
bracket is linear in $\xi$, the set of generators is the nullspace of a linear
operator. This observation drives a large body of work: it is how symmetries
of point clouds \citep{cahill2023liepca}, of trained networks
\citep{moskalev2022liegg,gruver2023lie} and of governing equations
\citep{kaiser2018discovering} are found, and it was recently unified and
substantially generalised by Otto et al.~\citep{otto2025unified}, who show that enforcing,
discovering and promoting symmetry are all governed by the same Lie-derivative
operator, with discovery reduced to computing a nullspace.

With real data the operator must be estimated, and an estimated operator has
no exact nullspace. Practice therefore replaces ``is this singular value
zero?'' with ``is this singular value small?'', by comparing it to a tolerance
or by looking for a gap in the spectrum. This fails in two ways.

\textbf{The tolerance means nothing in particular.} A singular value of the
Lie-derivative operator has the units of the data and depends on the arbitrary
scaling of the basis chosen for the candidate generators; a tolerance of
$10^{-3}$ is a statement about neither the system nor the evidence. Whether a
small value indicates a symmetry or merely a fortunate draw of the noise
depends on the noise level, the sample size, the conditioning of the design
and the dimension of the model class, none of which the tolerance sees.

\textbf{The burden of proof is inverted.} Suppose we improve matters by
attaching a hypothesis test, as a rank test would \citep{robin2000tests,
kleibergen2006generalized}: test $H_0$: ``this generator is an exact
symmetry'' and declare a symmetry when the test does not reject. Such a test
is valid, rejecting a true symmetry at most $5\%$ of the time, and it is still
the wrong instrument, because \emph{failing to reject $H_0$ is not evidence
for $H_0$}. The error we care about, declaring a symmetry that does not exist,
is the type-II error, and no level-$\alpha$ test controls it. In our
experiments this is not a technicality: a valid level-$0.05$ rank test
declares a symmetry that is not one in $7\%$ of replications on average over
our settings, and in every replication of its worst setting.

\paragraph{Our reframing.} A procedure that may honestly say ``$\xi$ is a
symmetry'' must be able to \emph{rule out} the alternative. So we ask for an
\emph{equivalence test}
\citep{schuirmann1987comparison,berger1996bioequivalence, wellek2010testing}:
fix a tolerance $\delta$, and certify $\xi$ only when the data reject the
hypothesis that its defect exceeds $\delta$. This buys error control in the
direction that matters: the probability of certifying a generator whose true
defect exceeds $\delta$ is at most $\alpha$. It also forces the tolerance into
the open, where it can be argued about.

Two ingredients make it practical. First, we measure symmetry violation with a
\emph{relative} defect: the norm of the bracket $[\xi,F]$ divided by the norm
of the two terms whose cancellation produces it (\Cref{def:defect}). The
resulting number lies in $[0,\sqrt 2]$, does not change when $\xi$ or $F$ is
rescaled or the candidate class is reparameterised, and answers a concrete
question: by what fraction do the two halves of the bracket fail to cancel?
Once the region of interest and the scales of the state variables are declared,
a tolerance in these units means the same thing for every system. Second, for
a model that is linear in its coefficients $\beta$ (a polynomial vector field,
or any fixed dictionary of functions), the squared defect is a \emph{ratio of
squared norms of linear images of $\beta$} (\Cref{lem:bilinear}). A confidence
ellipsoid for $\beta$ therefore propagates through it by optimising over an
ellipsoid, which reduces to trust-region subproblems
\citep{gander1989constrained,more1983computing} that can be solved globally and
with certified accuracy. Because a single event ($\beta$ lies in its
confidence region) implies the bound for \emph{every} generator at once, the
certificate is simultaneous over the whole candidate space, with no correction
for selection or multiplicity and no sample splitting. Each of these steps is
elementary on its own; the contribution is the formulation that makes them fit
together.

\paragraph{What we find.} We evaluate on ten polynomial dynamical systems
whose symmetry algebras we compute exactly by rational linear algebra, so that
every error is scored against ground truth rather than against a tolerance.
Across $240$ settings of system, sample size and noise level ($72{,}000$
replications), our
procedure, \symcert, never once certifies a generator whose true defect
exceeds the tolerance (the one-sided $95\%$ upper bound on its error rate is
$0.010$ in every setting). A fixed threshold at the same tolerance errs in
$0.5\%$ of replications on average and in $63\%$ in its worst setting; an
eigengap heuristic and the rank and split-sample significance tests reach
$100\%$, $100\%$ and $92\%$ in their worst settings. This error control costs
detection power, and we report the cost. Two non-polynomial systems fitted
with dictionaries that contain them behave in the same way. On real ocean
drifter velocities, the procedure certifies the rotation invariance that the
physics of near-inertial currents predicts, to a relative defect of $0.074$,
under assumptions stated with the result.

Four further findings, each a consequence of taking the statistics seriously.

\emph{Additive derivative noise rarely made a system without symmetry look
symmetric in our tests, but it can tilt a real symmetry.} When the right-hand side is measured with noise,
the class is correctly specified and the states are exact, the plug-in defect
of a direction near a true symmetry is biased \emph{upward}, which costs
detections. On a system without symmetry the smallest plug-in defect, selected
for being small, can fall below its true value, and a threshold at the
tolerance then declared a false symmetry in at most $3.7\%$ of replications,
and only when the smallest true defect was within a factor of about two of
the tolerance, never at four times it. The larger failure is on symmetric
systems: at high noise the estimated subspace tilts, and a fixed threshold
that declares the right number of generators can then declare some whose true
defect exceeds the tolerance; this is the source of the threshold rule's
$63\%$ worst case, on the planar rotation at the highest noise level.
\symcert{} certified no false symmetry in any of these tests.

\emph{Coverage has a computable price.} If the data occupy a region $\mu$
while the claim is about a domain $\nu$, the defect on $\nu$ can exceed the
defect measured on $\mu$ by at most an explicit factor built from two
generalised eigenvalues of Gram matrices (\Cref{prop:kappa}); the factor is
infinite exactly when the data lie in the zero set of a nonzero field of the
finite-dimensional space that contains every bracket, as happens, to numerical
precision, for a single trajectory that converges to an algebraic limit cycle
such as a circle; single trajectories of the Van der Pol and Lorenz systems
give large but finite factors. Across $50$ coverage settings, the nullspace
computed on the data declared spurious symmetry only beyond the threshold this
factor sets, and in none of the $13$ settings below it, where the proposition
rules out, at population level, declaring more generators than the true
algebra has; \symcert{} declared none in any setting.

\emph{The model class and noisy states can manufacture symmetry.} Fitting a
class too small to contain the truth produces a model with symmetries the
system does not have; sparse regression, the standard tool of data-driven model
discovery \citep{brunton2016discovering}, does this automatically on a weakly
symmetry-breaking system by deleting the term that breaks the symmetry; and
noise in the states makes least squares estimate a pseudo-true field, which on
a hardening Duffing oscillator was more symmetric than the system, enough for
the certificate itself to certify a false symmetry, while on a symmetric Hopf
system it hid the symmetry instead. Our certificate does not repair these
failures on its own, because its validity is conditional on a correctly
specified class and exactly observed states. A lack-of-fit test is not a
remedy either: we quantify how large a symmetry violation such a test of given
power can miss (\Cref{prop:lof}), and give a procedure that remains valid when the class is
chosen from the data, provided the largest class considered contains the truth
(\Cref{prop:classes}).

\emph{Units matter unless the metric is declared.} With a Euclidean norm on
the state, changing the units of one coordinate can turn a clear absence of
symmetry into an apparent one; declaring the scales of the state variables
removes the dependence (\Cref{lem:props}).

\paragraph{Contributions.}
\begin{itemize}\itemsep2pt
\item A relative equivariance defect (\Cref{def:defect}) that is invariant
  to the arbitrary choices in the nullspace formulation and, when the output
  metric is transported with the coordinates, to every affine change of state
  coordinates (\Cref{lem:props}); minimising it is a generalised eigenvalue
  problem.
\item \symcert: upper and lower confidence bounds on the defect of the
  \emph{true} system, finite-sample valid under Gaussian homoskedastic errors
  and as valid as the supplied coefficient region otherwise, simultaneously over all candidate
  generators and all tolerances, and hence over any data-dependent selection
  rule (\Cref{prop:simultaneous}). The bounds are the supremum and infimum of
  the defect over the confidence ellipsoid, not a ratio of separate bounds on
  its numerator and denominator, and are evaluated with certified numerics
  whose error can only widen them (\Cref{prop:closedform}). From the same region the procedure
  certifies approximate symmetry, rules out every approximate symmetry below
  a stated tolerance in the candidate class, and reports a certified plug-in
  subspace of approximate generators.
\item An explicit account of what the guarantee covers (\Cref{tab:guarantees}),
  including how large a symmetry violation a lack-of-fit test of given power
  can miss (\Cref{prop:lof}) and certificates that stay valid when the
  model class is chosen from the data within a nested family whose largest
  member contains the truth (\Cref{prop:classes}).
\item A non-asymptotic Le~Cam lower bound (\Cref{prop:lecam}) on the tolerance
  any procedure with false-certification rate $\alpha$ can certify; at a fixed
  model class, \symcert{} matches its $\sigma/\sqrt N$ scaling with a constant
  gap that we report, together with its growth in the number of coefficients.
\item A quantitative account of the mechanisms that create apparently valid
  symmetries: selection optimism, incomplete coverage, model truncation,
  sparse regression, noisy states and units, together with evidence that,
  in the systems tested, additive derivative noise made a system without
  symmetry look symmetric only rarely (in at most $3.7\%$ of replications)
  and only when its smallest true defect was within about a factor of two of
  the tolerance, although in a symmetric system it can tilt the subspace that
  a threshold declares.
\item An open-source implementation, exact ground truth for every polynomial
  benchmark, a full record of the compute, and a test suite.
\end{itemize}
